
\documentclass[a4paper,10pt]{report}

%\usepackage[landscape]{geometry}
\usepackage[cm-default]{fontspec}
\setromanfont{FreeSerif}
\setsansfont{FreeSans}
\setmonofont{FreeMono}
\usepackage[utf8x]{inputenc}
\usepackage{fontspec}
\usepackage{xunicode}
\usepackage{xltxtra}
\usepackage{xgreek}
\usepackage{amsmath}
\usepackage{unicode-math}
\usepackage{ulem}
\usepackage{color}
\usepackage{verbatim}
\usepackage{nopageno}
\usepackage{graphicx}
\usepackage{textpos}
\setlength{\TPHorizModule}{1cm}
\setlength{\TPVertModule}{1cm}
%\usepackage[colorgrid,texcoord]{eso-pic}
\usepackage[outline]{contour}
\usepackage{wrapfig}
\usepackage{url}
\usepackage{color}
\usepackage{tikz}
\usepackage{fancybox,fancyhdr}
\usepackage{subfigure}
\usepackage{pstricks}
\usepackage{epsfig}
\usepackage{multicol}
\usepackage{listings}
\usepackage{enumerate}
\usepackage{hyperref}
\hypersetup{
  bookmarks=true,
  bookmarksopen=true,
  pdfborder=false,
  pdfpagemode=UseNone,
  raiselinks=true,
  pdfhighlight={/P},
  colorlinks,
  citecolor=black,
  filecolor=black,
  linkcolor=black,
  urlcolor=black
}


\usepackage{fontspec}
\usepackage{xunicode}
\usepackage{xltxtra}


% Margins
%\setlength{\textwidth}{24cm}

%\setlength{\voffset}{0in}
%\setlength{\textheight}{6.8in}
% Colors
\definecolor{cbrown}{rgb}{0.49,0.24,0.07}
\definecolor{cmpez}{rgb}{0.92,0.88,0.79}
\definecolor{cmpez2}{rgb}{0.62,0.33,0.34}
\definecolor{cwhite}{rgb}{1,1,1}
\definecolor{cblack}{rgb}{0,0,0}
\definecolor{cred}{rgb}{0.9,0.15,0.15}
\definecolor{clblue}{rgb}{0.57,0.85,0.97}
\definecolor{corange}{rgb}{0.86,0.49,0.18}
\definecolor{clorange}{rgb}{0.95,0.84,0.61}
\definecolor{cyellow}{rgb}{1,0.95,0.16}
\definecolor{cgreen}{rgb}{0.69,0.8,0.31}
\definecolor{clbrown}{rgb}{0.78,0.67,0.43}
\definecolor{cmagenta}{rgb}{0.79,0.05,0.54}
\definecolor{cgray}{rgb}{0.76,0.73,0.63}
\definecolor{bluesite}{rgb}{0.05,0.43,0.67}
\definecolor{redsite}{rgb}{0.90,0.15,0.15}
\definecolor{backsite}{rgb}{0.07,0.07,0.07}
% ----------- from Costas -----------------------------

%---------------   Main Fonts    -----------------%   

% \setmainfont[Mapping=tex-text]{Calibri}
% \setmainfont[Mapping=tex-text]{Times New Roman}
%\setmainfont[Mapping=tex-text]{Droid Serif}
%\setmainfont[Mapping=tex-text]{Cambria}
%\setmainfont[Mapping=tex-text]{cm-unicode}
% \setmainfont[Mapping=tex-text]{Gentium}
% \setmainfont[Mapping=tex-text]{GFS Didot}
% \setmainfont[Mapping=tex-text]{Comic Sans MS}
 \setmainfont[Mapping=tex-text]{Ubuntu}
% \setmainfont[Mapping=tex-text]{Myriad Pro}
% \setmainfont[Mapping=tex-text]{CMU Concrete}
%\setmainfont[Mapping=tex-text]{DejaVu Sans}
%\setmainfont[Mapping=tex-text]{KerkisSans}
%\setmainfont[Mapping=tex-text]{KerkisCaligraphic}
% \setmainfont[Mapping=tex-text]{Segoe Print}
% \setmainfont[Mapping=tex-text]{Gabriola}

%------------------------------------------------%

% ------------ Mathematics Fonts ----------------%
\setmathfont{Asana-Math.ttf}
% -----------------------------------------------%
% Misc
\author{Κασωτάκη Ε. - Σμαραγδάκης Κ.}
\title{www.math24.gr}
% ---------------- Formation --------------------%


\setlength\topmargin{-0.7cm}
\addtolength\topmargin{-\headheight}
\addtolength\topmargin{-\headsep}
\setlength\textheight{26cm}
\setlength\oddsidemargin{-0.54cm}
\setlength\evensidemargin{-0.54cm}
%\setlength\marginparwidth{1.5in}
\setlength\textwidth{17cm}

\RequirePackage[avantgarde]{quotchap}
\renewcommand\chapterheadstartvskip{\vspace*{0\baselineskip}}
\RequirePackage[calcwidth]{titlesec}
\titleformat{\section}[hang]{\bfseries}
{\Large\thesection}{12pt}{\Large}[{\titlerule[0.9pt]}]
%--------------------------------------------------%
\usepackage{draftwatermark}
\SetWatermarkText{www.math24.gr}
\SetWatermarkLightness{0.9}
\SetWatermarkScale{0.6}
%--------------------------------------------------%


\usepackage{xcolor}
\usepackage{amsthm}
\usepackage{framed}
\usepackage{parskip}

\colorlet{shadecolor}{bluesite!20}

\newtheorem*{orismos}{Ορισμός}

\newenvironment{notation}
  {\begin{shaded}\begin{theorem}}
  {\end{theorem}\end{shaded}}




% Document begins
\begin{document}
\pagestyle{fancy}
\fancyhead{}
\fancyfoot{}
\renewcommand{\headrulewidth}{0pt}
\renewcommand{\footrulewidth}{0pt}


\fancyhead[LO,LE]{
 %\textblockcolor{backsite}
 %\begin{textblock}{5}(-2,-0.55)
  %\rule{0cm}{1cm}
 %\end{textblock}
 \textblockcolor{bluesite}
 \begin{textblock}{5}(-1.5,-0.55)
  \rule{0cm}{1cm}
 \end{textblock}
 %\textblockcolor{bluesite}
 %\begin{textblock}{14}(5.5,-0.55)
 % \rule{0cm}{1cm}
 %\end{textblock}
 \begin{textblock}{0}(-1,-0.25)
 \color{cwhite} \begin{Large}www.math24.gr\end{Large}
 \end{textblock}
\begin{textblock}{0}(16,-0.4)
 \color{cwhite} \includegraphics[height=1.5cm]{math24_logo.png}
 \end{textblock}
\textblockcolor{white}
\begin{textblock}{17}(-1,28)
 \color{backsite} \begin{small}Copyright \textcopyright 2012, Κασωτάκη Ε.(ikasotaki@gmail.com) - Σμαραγδάκης Κ.(kesmarag@gmail.com)\end{small}
 \end{textblock}
}
 
\begin{shaded}
\begin{center}
\huge \textbf{Φυλλάδιο Ασκήσεων}\\
%Πρόσθεση ρητών αριθμών
\end{center} 
\textbf{Μαθηματικά Γ' Γυμνασίου} \hfill \textbf{Ημερομηνία Παράδοσης : \hspace{2em} }
\subsection*{Ονοματεπώνυμο :\hfill  \hspace{5em}}
\end{shaded}
\vspace{2em}
\begin{itemize}
 \item Αριθμητικές και αλγεβρικές παραστάσεις
 \item Ακέραια αλγεβρική παράσταση
 \item Αριθμητική τιμή
 \item Μονώνυμα
 \item Συντελεστής, κύριο μέρος και βαθμός μονωνύμων
 \item Όμοια, ίσα και αντίθετα μονώνυμα
 \item Σταθερό και μηδενικό μονώνυμο
\end{itemize}

\section*{Θεωρία  \hfill \small{}}
\begin{itemize}
\item \textbf{Αριθμητικές παραστάσεις} ονομάζονται οι εκφράσεις που περιέχουν μόνο αριθμούς 
       (αριθμούς που συνδέονται με πράξεις).\\
       \textbf{Π.χ} $4+3$ \\
       \textbf{Π.χ} $12\cdot 4 -3^{2}+7\cdot 8$ 
\item \textbf{Αλγεβρικές παραστάσεις} ονομάζονται οι εκφράσεις που περιέχουν αριθμούς και μεταβλητές 
       (που συνδέονται με πράξεις).\\
       \textbf{Π.χ} $4xy^{2}$ \\
       \textbf{Π.χ} $8x+9y^{2}+xy$ \\
       \textbf{Π.χ} $\dfrac{3x^{2}}{4y}$ 
\end{itemize}


\section*{Θεωρία  \hfill \small{}}
\begin{itemize}
 \item \textbf{Ακέραια} αλγεβρική παράσταση ονομάζεται μία αλγεβρική παράσταση όταν μεταξύ των 
        μεταβλητών της σημειώνονται μόνο οι πράξεις της πρόσθεσης και του πολλαπλασιασμού και 
        οι εκθέτες των μεταβλητών της είναι φυσικοί αριθμοί.\\
        \textbf{Π.χ} $4x^{3}+5x$
 \item \textbf{Αριθμητική τιμή} ή \textbf{τιμή} μιας αλγεβρικής παράστασης ονομάζεται ο αριθμός που 
        θα προκύψει εάν αντικαταστήσουμε τις μεταβλητές της αλγεβρικής παράστασης με αριθμούς και 
        κάνουμε τις πράξεις.\\
        \textbf{Π.χ} Η αλγεβρική παράσταση $x^{2}+3x$ για $x=2$ έχει τιμή $10$ γιατί: $2^{2}=3\cdot 2=4+6=10$
\end{itemize}



\section*{Άσκηση 1  \hfill \small{ μονάδες}}
Να βρείτε την αλγεβρική τιμή των παρακάτω αλγεβρικών παραστάσεων:
\begin{enumerate}[i)]
 \item $4x^{3}-2x+17y$ για $x=2$ και $y=1$
 \item $3x^{2}y-2xy^{2}$ για $x=-1$ και $y=3$ 
 \item $12αβ^{2}+18α^{2}β$ για $α=\dfrac{1}{3}$  και $β=\dfrac{1}{2}$
 \item $2\sqrt{x}+xy^{2}-\dfrac{1}{2}x^{2}y$ για $x=4$ και $y=5$
\end{enumerate}

\section*{Θεωρία  \hfill \small{}}
\begin{itemize}
 \item \textbf{Μονώνυμα} ονομάζονται οι ακέραιες αλγεβρικές παραστάσεις στις οποίες μεταξύ των μεταβλητών σημειώνεται 
       μόνο η πράξη του πολλαπλασιασμού.\\
       \textbf{Π.χ} $\quad 5x^{3}$ \\
       \textbf{Π.χ} $\quad 7y^{2}$ \\
       \textbf{Π.χ} $\quad 2x^{4}y^{5}$
 \item \textbf{Συντελεστής} του μονωνύμου ονομάζεται ο αριθμητικός παράγοντας του. \\
       \textbf{Π.χ} το μονώνυμο $6x^{2}$ έχει συντελεστή το $6$
 \item \textbf{Κύριο μέρος} του μονωνύμου ονομάζεται το γινόμενο όλων των μεταβλητών με τους 
       αντίστοιχους εκθέτες τους.\\
       \textbf{Π.χ} το μονώνυμο $7xy^{2}$ έχει κύριο μέρος το $xy^{2}$
 \item \textbf{Βαθμός} του μονωνύμου 
        \begin{itemize}
         \item  \textbf{ως προς μία μεταβλητή} είναι ο εκθέτης της μεταβλητής αυτής
         \item  \textbf{ως προς όλες τις μεταβλητές} είναι το άθροισμα των εκθετών των μεταβλητών του 
        \end{itemize}
       \textbf{Π.χ} το μονώνυμο $4x^{2}y^{5}$ 
       \begin{itemize}
        \item είναι δευτέρου βαθμού ως προς τη μεταβλητή $x$ 
        \item είναι πέμπτου βαθμού ως προς τη μεταβλητή $y$ 
        \item είναι έβδομου βαθμού ως προς τις μεταβλητές $x$ και  $y$ 
        \end{itemize}
\end{itemize}


\section*{Άσκηση 2  \hfill \small{ μονάδες}}
Να συμπληρώσετε τον παρακάτω πίνακα επιλέγοντας \textbf{"X"} στην κατάλληλη θέση:
\begin{center}
 \begin{tabular}{|c|c|c|}
\hline                   
        \textbf{Αλγεβρική παράσταση} &  \textbf{Είναι μονώνυμο} & \textbf{Δεν είναι μονώνυμο}     \\
\hline 
$2x+5x^{2}$                          &                          &        \\
\hline
$6x^{3}y$                            &                          &   \\
\hline
$\dfrac{7}{2}xy^{2}$                 &                          &   \\
\hline
$3\dfrac{x^{2}}{y}$                  &                          &   \\
\hline
$y^{3}x^{5}$                         &                           &  \\
\hline
\end{tabular}
\end{center}








\section*{Άσκηση 3  \hfill \small{ μονάδες}}
Να συμπληρώσετε τον παρακάτω πίνακα: 
\begin{center}
 \begin{tabular}{|c|c|c|c|c|c|}
\hline                   
 \textbf{Μονώνυμο}& \textbf{Συντελεστής} &\textbf{Κύριο μέρος} & \textbf{Βαθμός ως προς $x$}  &  \textbf{Βαθμός ως προς $y$}  & \textbf{Βαθμός ως προς $x$ και $y$}    \\
\hline 
$3x^{2}y^{3}$   &   &   & &  & \\
\hline
$4xy^{2}$   &   &   &   &  & \\
\hline
$\dfrac{2}{7}$ &  & &  &  & \\
\hline
$-\sqrt{5}x^{5}$ &  & & & & \\
\hline
$-xy^{4}$  & & & & & \\
\hline
\end{tabular}
\end{center}

\section*{Θεωρία  \hfill \small{}}
\begin{itemize}
 \item \textbf{Όμοια} ονομάζονται τα μονώνυμα που έχουν το ίδιο κύριο μέρος.\\
       \textbf{Π.χ} \quad $3x^{2}y$ και $-7x^{2}y$ είναι όμοια        
 \item \textbf{Ίσα} ονομάζονται τα όμοια μονώνυμα αν έχουν τον ίδιο συντελεστή.\\
       \textbf{Π.χ} \quad $4xy^{2}$  και $\dfrac{8}{2}xy^{2}$ είναι ίσα
 \item \textbf{Αντίθετα} ονομάζονται τα όμοια μονώνυμα που έχουν αντίθετους συντελεστές.\\
       \textbf{Π.χ} \quad $6xy^{3}$ και $-6xy^{3}$ είναι αντίθετα
\end{itemize}


\section*{Άσκηση 4  \hfill \small{ μονάδες}}
Να αντιστοιχίσετε κάθε μονώνυμο της αριστερής στήλης με ένα όμοιο μονώνυμο της δεξιάς στήλης
\begin{itemize}
\begin{multicols}{2}
 \item $3xy^{2}$
 \item $-\sqrt{2}x^{3}y^{3}$
 \item $-4x^{5}$
 \item $2x^{2}y^{2}$
 \item $3x^{2}y$
%---------------------------------
 \item $-9x^{5}$
 \item $5xy^{2}$
 \item $-9x^{2}y^{2}$
 \item $10x^{3}y^{3}$
 \item $\sqrt{5}x^{2}y$
\end{multicols}
\end{itemize}


\section*{Θεωρία  \hfill \small{}}
\begin{itemize}
 \item \textbf{Σταθερό} μονώνυμο είναι οι αριθμοί (γιατί τους θεωρούμε ως μονώνυμα) και είναι μηδενικού βαθμού.
 \item \textbf{Μηδενικό} μονώνυμο θεωρείται ο αριθμός $0$ και δεν έχει βαθμό.
\end{itemize}


\section*{Άσκηση 5  \hfill \small{}  μονάδες}
Να βρείτε ένα μονώνυμο
\begin{enumerate}[i)]
 \item που έχει συντελεστή $6$ και κύριο μέρος $x^{5}y^{3}$
 \item που είναι ίσο με το $7xy^{2}$
 \item που είναι αντίθετο με το $-13x^{3}y^{2}$
 \item που έχει συντελεστή $\sqrt{2}$ και είναι όμοιο με το $5x^{4}y^{3}$ 
 \item που έχει συντελεστή $2$ και μεταβλητές $x$ και $y$, και βαθμό ως προς $x$ είναι $3$ και ως 
       προς $y$ είναι $4$
 \item που έχει συντελεστή $-7$ και μεταβλητές $x$ και $y$, και ο βαθμός του ως προς $x$ είναι $4$ 
       και ο βαθμός του ως προς $x$ και $y$ είναι $5$
\end{enumerate}



\section*{Άσκηση 6  \hfill \small{}  μονάδες}
Να προσδιορίσετε την τιμή του φυσικού αριθμού $ν$ ώστε το μονώνυμο $5x^{2}y^{ν}$ 
\begin{enumerate}[i)]
 \item να είναι τρίτου βαθμού ως προς $y$
 \item να είναι τρίτου βαθμού ως προς $x$ και $y$
\end{enumerate}

 









%\let\thefootnote\relax\footnotetext{Επιμέλεια : Κασωτάκη Ειρήνη (ikasotaki@gmail.com),    Σμαραγδάκης Κώστας (kesmarag@gmail.com)} 


\end{document}
